\section{End-to-End Evaluation}
\label{sec:endtoend}

% CENTERPIECE of the system paper. The deployed system, evaluated top-down:
% lead with the headline (deflection fix), then the full 3x3 matrix that
% justifies the deployed config, then the safety/usefulness tradeoff and the
% capability ceiling. Numbers: pull exact values from mamai-eval
% configs/config-v0.2.0/reports/generator-prompt-matrix-20260622.html and
% the dangerous-case adjudication (dangerous-case-adjudication-20260623.md).
% The webpage rounds; the eval reports are the source of record.

\subsection{Headline: Fixing Deflection}
% The base model deflected on ~1/3 of questions ("see a doctor"). The G1
% prompt (Section ref sec:system-prompt) cut deflection 32.7% -> 3.2% and
% ~doubled key-fact recall (0.139 -> 0.279) with no added dangerous answers.
% This is the single largest end-to-end gain and the reason G1 ships.

\subsection{The Generator x Prompt Matrix}
% Full 3x3: generators {Gemma 4 E4B (deployed), Gemma 3n E4B, Qwen3.5-397B
% (frontier ceiling)} x prompts {baseline, +G1, +G1+G2}. Scored on two
% open-ended benchmarks from mamabench: Kenya Clinical Vignettes (n=312,
% nurse-written, close to deployment) and HealthBench-oss (n=1,209,
% physician-rubric). Metrics per cell: key-fact recall, deflection,
% potentially-harmful %, dangerous count (Kenya); weighted_met with the
% positive-credit vs penalty decomposition (HealthBench). Identical RAG
% pipeline across all cells (EmbeddingGemma top-3, bundle v0.3.0); only the
% generator/prompt changes. Judge: gpt-oss-120b (validated, Section
% ref sec:methodology). Present the 9-cell table here.

\subsection{Safety versus Usefulness}
% The two on-device models trade off on a single axis. Gemma 3n is more
% useful (recall 0.30-0.39, ~2x HealthBench) but less safe: 4 -> 11 -> 15
% dangerous answers across prompts. Gemma 4 is the mirror: less complete but
% near-zero dangerous (0 -> 1 -> 0). G1/G2 raise engagement, which on a weak
% grounder (3n) concentrates harm exactly where it is weakest (dosing).
% Decision: deploy Gemma 4 (safe), use G1 to win back most of the usefulness.

\subsection{Dangerous-Case Adjudication}
% All 31 dangerous-flagged cases hand-read (dangerous-case-adjudication-
% 20260623.md). Gemma 3n's are genuine serious errors: oxytocin dosed in mg
% not milli-units (~1000x overdose), ceftriaxone in a neonate
% (contraindicated), benzathine penicillin ~10x underdose, PEP eligibility
% reversed. Gemma 4's single +G1 flag is a judge over-flag (standard choking
% first-aid). This is what grounds the "safe" claim in clinical terms.

\subsection{The Capability Ceiling}
% Qwen3.5-397B (unconstrained frontier) shows the prompts are sound: G1+G2 is
% pure upside (recall 0.553, harm 4.2%, 0 dangerous). The deployed config
% still trails this ceiling on recall/completeness/safety alike. The gap is
% generator capability and grounding, not prompt design -- motivating the
% component layers that follow and the fine-tuning direction (Section
% ref sec:discussion).

\subsection{MCQ as a Negative Control}
% Aside: full MCQ suite +-RAG on Gemma 4 (n=23,241): no-RAG 53.7% vs +RAG
% 51.9% (-1.8 pp). Not a deployment failure -- the maternal/neonatal corpus
% is off-domain for broad USMLE-style MCQ, so retrieved chunks distract.
% This is why the deployment-relevant signal is the open-ended tracks above,
% not MCQ. (mcq-rag-effect-20260520.md)
